\section{Introduction}

Redundant manipulators can exploit extra degrees of freedom for tracking, posture regulation, joint limit avoidance, and singularity avoidance. In repetitive tasks, however, redundancy also creates the classical joint drift problem: the end effector may return to a periodic Cartesian path while the joints fail to recover the initial posture at the end of each cycle. Once this cycle end error accumulates, the manipulator can gradually lose manipulability or move toward joint limits. For this reason, drift free repetitive motion planning remains a representative benchmark problem in robot kinematic control \citep{zhang2009rmp,zhang2017three}.

The classical route is to formulate repetitive inverse kinematics as a constrained quadratic program and to solve it online by recurrent neural dynamics \citep{wang2004recurrent,toshani2014ik,tan2022physical,hassan2020review}. In the drift free setting of \citet{zhang2009rmp}, the objective combines posture recovery, task equality, and dynamic velocity bounds converted from joint limits. This formulation is elegant and still influential, but its derivation and solver interpretation come mainly from a continuous time neural dynamics viewpoint. In practice, industrial and simulator based robot control is executed in sampled form, with explicit control periods, finite inner iterations, and communication induced mismatch. The sampled controller therefore needs a discrete interpretation rather than a direct inheritance from continuous time analysis.

LCCZNN and DLCCZNN offer an appealing route because they replace dense matrix inversion by residual driven iterative dynamics and preserve a clear error system interpretation \citep{zheng2024lccznn,qiu2022jerk,yang2025predefined,wen2024firstsecond,li2026dlccznn}. Yet the present comparative study also shows that direct solver replacement is not automatically beneficial. Some replacements preserve the drift free closed loop property, whereas others produce small instantaneous position errors but destroy cycle end posture recovery. This indicates that the real control issue is not only numerical efficiency, but tight coupling between algorithmic dynamics and control structure under a finite sampling budget.

The paper is built around a question that is easy to state but often blurred in existing work: when moving from continuous time drift free planning to sampled robot control, should one change the objective, change the solver, or do both? This paper answers that question in a deliberate order. The original \methodone{}, \methodtwo{}, and \methodthree{} are first retained as baseline formulations. Among them, \methodthree{} reaches the highest raw accuracy because it changes the objective itself, whereas \methodtwo{} is the strongest baseline that still preserves the classic drift free kinematic resolution formulation. This makes \methodtwo{} the right anchor for the main contribution of the paper: a sampled \methodtwoDL{} solver derived directly from the code level implementation and interpreted in control terms. To avoid an overly narrow story, three auxiliary lines are also retained: \methodoneDL{} as a weaker formulation preserving transfer, \methodthreelcc{} as a direct LCCZNN replacement of the Method~3 objective, and \methodthreepcg{} as a secondary structured solver for the Method~3 discrete subproblem.

The main contributions are as follows.
\begin{itemize}[leftmargin=1.5em]
  \item A formulation preserving sampled \methodtwoDL{} controller is proposed for the classic drift free kinematic resolution problem. Starting from the strongest classical formulation preserving baseline, its update law tracks the sampled Karush Kuhn Tucker residual by projected DLCCZNN dynamics, so the solver remains close to the logic of sampled robot control rather than continuous time transcription.
  \item A discrete control interpretation is established for the classical drift free repetitive motion quadratic program and for the proposed controller, including the explicit role of task feedback velocity, one step error propagation, dynamic velocity bounds induced by physical limits, and the sampled residual to command mismatch relation.
  \item A unified comparison is built across the original baseline methods and the auxiliary solver variants \methodoneDL{}, \methodthreelcc{}, and \methodthreepcg{}. This comparison reveals which gains come from formulation preserving solver design and which gains come from changing the optimization objective itself, thereby giving the paper a clear scientific hierarchy instead of a flat list of methods.
  \item Offline UR3e simulations and simulator based real time control loop validation in CoppeliaSim show that the proposed \methodtwoDL{} improves the classical \methodtwo{} under sampled execution, while the auxiliary comparisons clarify both the promise and the limits of alternative solver transfers. The parameter study further shows that offline and real time optima do not coincide, so drift recovery under sampled execution must guide the final retained setting.
\end{itemize}

The remainder of the paper is organized as follows. Section~2 reviews drift free repetitive motion planning and low complexity neural solvers. Section~3 formulates the discrete control problem. Section~4 presents the baseline methods, the proposed \methodtwoDL{}, and the auxiliary solver variants. Section~5 analyzes the proposed solver. Sections~6 and 7 report the experimental design and results. Section~8 concludes the paper.
